%%%%%%%%%%%%%%%%%%%%%%%%%%%%%%%%%%%%%%%%%
% Laboratory report template
% 
% Template made for the course Basic Chemistry
% at Stockholm University
%
% Template by Viktor Bengtsson
%
% Template creation date: 2018-08-24
%
%%%%%%%%%%%%%%%%%%%%%%%%%%%%%%%%%%%%%%%%%

%----------------------------------------------------------------------------------------
%	PACKAGES AND DOCUMENT CONFIGURATIONS
%----------------------------------------------------------------------------------------

\documentclass[11pt,a4paper,british,swedish]{article}
\usepackage[T1]{fontenc}
\usepackage[utf8]{inputenc}
\usepackage{times} % Use Times Font

\usepackage{booktabs} % Allow nicer tables
\usepackage[pdftex]{graphicx} % Required for including pictures
\usepackage[space]{grffile} % Allows using spaces in filenames for images
\usepackage[british,swedish]{babel} % Swedish translations

\usepackage[pdftex,hidelinks]{hyperref} % Format links for pdf
\usepackage[noabbrev]{cleveref} % Easier to make automated references
\usepackage{bookmark} % Format bookmarks

\usepackage{calc} % To reset the counter in the document after title page
\usepackage{enumitem} % Includes lists
\usepackage{amsmath} % Required for some math elements 
\usepackage[version=4]{mhchem} % Easily write chemical formulas
\usepackage{siunitx} % Easily format SI units
\sisetup{output-decimal-marker = {,}} % Makes sure that \si{} and \SI{} outputs decimals with comma

\usepackage{datetime2} % Chages date format of \today to the suitable one

\usepackage[safeinputenc,style=science,backend=biber]{biblatex} % Referencence handler, making references behave as for "Science" publications
\usepackage{csquotes} % Quotation formatting

\frenchspacing % No double spacing between sentences
\linespread{1.2} % Set linespace
\usepackage[a4paper, lmargin=0.1666\paperwidth, rmargin=0.1666\paperwidth, tmargin=0.1111\paperheight, bmargin=0.1111\paperheight]{geometry} %margins
%\usepackage{parskip}
\usepackage{fancyhdr} % Allows inserting headers

\usepackage[all]{nowidow} % Tries to remove widows
\usepackage[protrusion=true,expansion=true]{microtype} % Improves typography, load after fontpackage is selected

\usepackage{lipsum} % Used for inserting dummy 'Lorem ipsum' text into the template


%----------------------------------------------------------------------------------------
%	DOCUMENT INFORMATION
%----------------------------------------------------------------------------------------

\addbibresource{Zotero.bib} % Add bibliography resource. Current name to be connected with zotero in Overleaf

% Lägg till information som ska vara med på flera ställen
\newcommand{\laboration}{Labb X: Lorem Ipsum} % laborationens namn, exempelvis Syntes
\newcommand{\laborant}{Laborantens Namn}      % laborantens namn, exempelvis John Doe
\newcommand{\kursnamn}{Grundläggande kemi}


% Lägg till information till titelsidan (görs automatiskt utifrån kommandon ovan)
\title{Mall för labbrapporter} % Lägg till titel
\author{\laborant}
\date{\today} % Lägger till dagens datum när dokumentet skapas. \today kan ersättas med ett annat datum.


% Definiera sidhuvud
\pagestyle{fancy}
\fancyhf{}
\setlength{\headheight}{14pt}
\fancyhead[L]{\kursnamn}
\fancyhead[C]{\laboration}
\fancyhead[R]{\laborant}
\fancyfoot[C]{\thepage}


%-----------------------
% Begin document
%-----------------------
\begin{document} % All text i dokumentet hamnar mellan dessa taggar, allt ovanför är formatering av dokumentet
\maketitle % Gör så att titel, namn och datum läggs in i dokumentet
% Please add the following required packages to your document preamble:
% \usepackage{graphicx}
\begin{center}
\begin{tabular}{lr}
Kurs:                         & \kursnamn                        \\
Kursdel:                      & Jämvikt                          \\
Laboration:                   & \laboration                      \\
Medlaborant:                  & Medlaborantens Namn              \\
Labbhandledare:               & Viktor Bengtsson                 \\
Utförandedatum:               & \DTMdisplaydate{2018}{9}{5}{-1}

\end{tabular}%
\end{center}

\newpage % Tvingar en sidbrytning


\section{Inledning\label{sec:inledning}} % \label{sec:inledning} gör så du kan referera hit med
% \ref{sec:inledning}. sec syftar på "section". Figurer hålls reda på med "fig:" och tabeller med
% "tab:". Det går att skriva utan de namnen också, men då blir det svårare att veta vilka etiketter
% som hör till vad.

% Här skriver du inledning med lämplig mängd teori, syfte och liknande

\subsection{Underrubrik} % Här kan du lägga till en underrubrik om du vill ha den. Annars tar du bort raden
% Exempel på reaktionsformel:
\begin{center}
\ce{Na2SO4 -> 2Na+(aq) + SO4^2-(aq)}
\end{center}

% Enheter:
\si{\mol\per\cubic\deci\meter}

\SI{1,0}{\mol\per\cubic\deci\meter}


\section{Metod\label{sec:metod}}




\section{Resultat\label{sec:resultat}}
% Tabeller görs förslagsvis via https://www.tablesgenerator.com/latex_tables
% https://www.latex-tables.com eller liknande och kopieras in

\begin{figure}
	\centering
	\includegraphics[width=\linewidth]{solar-spectrum} % Infoga namnet på filen innanför {}. Den måste finnas bland "files"
	\caption{Figurtexten skrivs här \label{fig:}}
\end{figure}



\section{Diskussion och slutsats\label{sec:diskussion}}


\printbibliography
\end{document}
